\documentclass{report}
\usepackage{amsmath,amssymb}
\setlength{\parindent}{0mm}
\setlength{\parskip}{1em}
\begin{document}
\begin{center}
\rule{6in}{1pt} \
{\large Paul G. Constantine \\
{\bf A Lanczos Method for Approximating Composite Functions}}

2286 Williams Street \\ Palo Alto \\ CA 94306
\\
{\tt paul.constantine@stanford.edu}\\
Eric T. Phipps\end{center}

Many complex multiphysics models employ composite functions, where each
member function represents a different physics. A simple example of this
is a chemical reaction model; the decay of the concentration depends on
the decay rate parameter, but the model for the decay rate (e.g., the
Arrhenius model) depends on the temperature, the gas constant, the
activation energy, and the prefactor. We consider the general setting
\[
f = f(x) \quad g = g(f)
\]
where $x$ are the inputs to the physics defined by $f$, and $f$ are the
inputs for the physics defined by $g$. One may be interested in
understanding how $g$ behaves as $x$ changes, so sensitivity and
uncertainty studies can be performed on the composite function $h(x) =
g(f(x))$. Evaluating $h$ is often a computationally demanding task,
rendering studies that require many evaluations infeasible, particularly
when the dimension of $x$ is large. In this work, we propose a strategy
to take advantage of the composite structure of $h$ to build surrogate
models. The strategy is particularly advantageous when $g(f)$ is much
more expensive to evaluate than $f(x)$. The essence of the strategy is to
use the relatively cheap evaluations of $f$ to determine a small set of
points in its range space to evaluate $g$. This is especially applicable
when the dimension of $x$ is large -- when methods that require
evaluating $h$ at many points in the high dimensional $x$-space become
infeasible.

The strategy is closely linked to Gaussian quadrature. We implicitly
approximate the density function of $f$ and construct a set of
polynomials of $f$ that are orthonormal with respect to its density
function. The function $g$ is then approximated as a truncated series in
these basis polynomials of $f$, which is in contrast to the standard
methods of approximating $h$ as an orthonormal polynomial series in $x$.
In the context of uncertainty quantification, such polynomial
approximations appear under the names polynomial chaos or stochastic
collocation, amongst others. We use a discrete Stieltjes procedure to
compute the recurrence coefficients of the orthogonal polynomials in $f$,
and we show how this is equivalent to a Lanczos' method on a diagonal
matrix with a weighted inner product. The basis vectors from the Lanczos
iteration can be used to linearly map a few evaluations of $g(f)$ to many
evaluations of $g(f(x))$, which can then be used to study dependence of
$h$ on $x$.


\end{document}
